\section{Support of the Rees pieces and the upper valuation estimate}
\label{sec:deformation-support}

We prove the last inequality in \eqref{eq:single-rees} from the
graded presentation and its point conditions. Write \(D=A_0\) and
\(I_0\) for the homogeneous ideal of the original equations. Use
\(L_0,L_1,L_2,\mathcal R,s_*\) from the preceding two sections and
write \(Q=\mathcal R/J\), where \(J\) is any homogeneous ideal
containing all the sections \(s_*\) and all the degree-zero \(Z_i\).
The module to which the upper estimate applies is \(P_d=Q_d\).
In the application \(J\) also contains the coefficient kernel of
the chart map. We need only the inclusions just stated.

We recall the assumptions about points of the presentation used here.
The ring \(D\) is Noetherian local, with the same residue field as
\(\mathcal O\), and \(D/(D\cap I_0)\) is finite as an
\(\mathcal O\)-module. The weights of a finite set of generators
for \(A\) over \(D\) lie between \(-2\) and \(2\). The weight
\(-1,-2\) equations span their respective saturated cotangent
lattices after extending scalars to \(F\). At a fixed geometric
point on the equations other than the augmentation, the model
excludes simultaneous existence of two solutions with the same
values on \(D\), one supported on nonpositive weights but nonzero
on some negative weight, and one supported on nonnegative weights
but nonzero on some positive weight. This is the specific geometric
assertion called \(M3\). The mixed-point tangent condition used in
Section~\ref{sec:deformation-chart} is also assumed.

\subsection{Equations seen by a relevant prime}

Let \(U_i\) be the ideal of \(\mathcal R\) generated by \(L_i T_i\),
and put \(\mathfrak b=U_0U_1U_2\). Call a prime \(\mathfrak P\)
of \(\mathcal R\) relevant if it contains \(J\) and does not contain
\(\mathfrak b\). For such a prime choose
\[
 f\in A_{-2},\quad g\in A_2,\quad h\in L_2,\qquad
 fT_0,\ gT_1,\ hT_2\notin\mathfrak P .
\]
The weights of \(f,g,h\) are \(-2,2,-2\). Write
\(\mathfrak p=\mathfrak P\cap D\). We show
\begin{align}
 D\cap I_0&\subseteq\mathfrak p, \label{eq:rees-contract}\\
 \Delta\subseteq\mathfrak p&\ \Longrightarrow\
         \mathfrak p=\mathfrak m_D,\qquad
 \Delta=(A_{-2}A_2)D\ \subseteq D . \label{eq:rees-boundary}
\end{align}
In the formula for \(\Delta\) the ideal is generated by the indicated
products.

First suppose \(\tau=uv\notin\mathfrak p\) for \(u\in A_{-2}\),
\(v\in A_2\). Localize at \(\tau\), so \(u,v\) are units. The
weight \(2j\) component of the localized algebra is \(u^{-j}D_\tau\)
by multiplication by the unit \(u^j\). In particular the three
modules \(L_0,L_1,L_2\) become \(uD_\tau,u^{-1}D_\tau,uD_\tau\);
the last equality uses \(A_{-2}^2 A_2\subseteq L_2\subseteq A_{-2}\).
Thus the localized Rees algebra is the polynomial algebra over
\(D_\tau\) on \(uT_0,u^{-1}T_1,uT_2\). Relevance says none of
these variables belongs to the localized prime. The sections force
the coefficient of every corresponding weight-normalized generator
\(F_i^2,G_i,E_i^2,E'_i,Z_i\) into its contraction. For a homogeneous
\(a\) of weight opposite to an unsquared equation, \(a\) times that
equation has degree zero. Its square is a scalar multiple of the
corresponding normalized squared equation in the odd-weight cases,
so the product belongs to the same prime. In even weights one uses
the normalized unsquared equation directly. Every element of
\(D\cap I_0\) is a sum of such degree-zero products by homogeneity.
This proves \eqref{eq:rees-contract} on this open set.

Suppose next that \(\Delta\subseteq\mathfrak p\). We construct
one-sided points directly. For either sign \(\epsilon=\pm1\), put
\(t=g,i=1\) if \(\epsilon=1\), and \(t=f,i=0\) if \(\epsilon=-1\).
The even components satisfy
\(A_{2\epsilon n}=(A_{2\epsilon})^n\) for \(n\ge0\), with products
over \(D\). Indeed in any generator monomial of weight \(2\epsilon n\)
the factors of sign \(\epsilon\) have total absolute weight at least
\(2n\). One can choose \(n\) disjoint blocks of weight \(2\epsilon\),
each a single weight-two factor or two weight-one factors. All
remaining factors have net weight zero.

In the fraction field of \(\mathcal R/\mathfrak P\), define on
homogeneous \(a\in A_{2\epsilon m}\), \(m\in\mathbb Z\),
\[
 \theta_\epsilon(a)=
       \frac{[(at^k)T_i^{\,m+k}]}{[tT_i]^k},
             \qquad k\ge\max(0,-m).
\]
The even-component identity makes the numerator a Rees element.
Increasing \(k\) multiplies numerator and denominator by the same
nonzero element. Choosing common \(k\)'s proves additivity, and
using the sum of the two chosen \(k\)'s proves multiplicativity.
Thus \(\theta_\epsilon\) is a ring map on the even Veronese subring.
It induces the original map on \(D\), and its value on \(t\) is
nonzero. For \(a\in A_{-2\epsilon}\), the numerator with \(k=1\)
is \(at\), which belongs to \(\Delta\), so its value is zero.
The even-component identity kills every weight of opposite even sign.

All squared odd-weight equations on the chosen side, and the
unsquared equations of weight \(2\epsilon\), map to zero since their
sections belong to \(\mathfrak P\). For the negative side this also
holds for \(F_{r+1}^2T_0\), which was imposed in slot two. Indeed,
\[
        (F_{r+1}^2T_0)(hT_2)=(hT_0)(F_{r+1}^2T_2)
\]
and \(hT_2\notin\mathfrak P\). The degree-zero equations map to
zero as well. The algebra \(A\) is integral over its even subring:
an odd element squares into that subring, and every element is a
sum of an even and an odd element. By lying over, extend the map
to \(A\) with values in a field extension. Odd equations on the
chosen side vanish because their squares do. Every element of odd
opposite sign vanishes for the same reason. We have obtained a
solution of the equations on which only the indicated sign and
weight zero can be nonzero, with \(t\) surviving. Its restriction
to \(D\) has kernel \(\mathfrak p\). In particular this proves
\eqref{eq:rees-contract} also on the boundary.

It remains to exclude a boundary prime in characteristic zero.
By \eqref{eq:rees-contract}, \(D/\mathfrak p\) is finite over
\(\mathcal O\). If \(\mathfrak p\) meets \(\mathcal O\) nontrivially,
it contains \(\varpi\), so this quotient is a finite-dimensional
domain over the residue field and is a field. Thus
\(\mathfrak p=\mathfrak m_D\). Otherwise choose an embedding
of \(\operatorname{Frac}(D/\mathfrak p)\) in \(k=\overline F\).
Both one-sided solutions can be realized in \(k\) over this fixed
embedding. For each sign the equations, the requirement that
opposite generators vanish, and \(t\ne0\) define a finite type
algebra over \(\operatorname{Frac}(D/\mathfrak p)\). The constructed
field-valued solution shows it is nonzero. Extending scalars to
\(k\) preserves nonzeroness, and a maximal ideal gives the desired
\(k\)-point.

Sending every nonzero-weight element of the solution of negative
sign to zero and preserving its values on \(D\) produces a fixed
point. One way to see multiplicativity is to send a homogeneous
element of weight \(j\le0\) to its value times \(T^{-j}\) and send
all \(j>0\) to zero. The original solution vanishes on positive
weights, so this gives a map to \(k[T]\) whose specialization at
zero is the fixed point. The fixed point still satisfies the
homogeneous equations. If it differs from the augmentation, the
two one-sided solutions contradict \(M3\).

If it is the augmentation, the one-sided solution \(\xi_-\) agreeing
with it on \(D\) must vanish also on negative weights. Here are the
details. A homogeneous element of \((\ker x_A)^2\) of weight \(-1\)
is a sum of products of two homogeneous augmentation elements whose
weights sum to \(-1\). At least one factor has nonnegative weight,
and its value at \(\xi_-\) is zero. The same holds for the saturation
since nonzero scalars remain nonzero. Spanning by the \(F_i\) over
\(F\) now makes \(\xi_-\) vanish on all of \(A_{-1}\). In weight
\(-2\), products of two augmentation elements either have a
nonnegative-weight factor or have two factors of weight \(-1\).
So the saturated square vanishes there as well, and spanning by the
\(G_i\) kills \(A_{-2}\). This contradicts \(\xi_-(f)\ne0\).
This proves \eqref{eq:rees-boundary}.

\subsection{Passing to large multidegrees}

Put \(I_D=D\cap I_0\). For every prime of \(\mathcal R\) containing
\(J\), either \(\mathfrak b\) or \(I_D\) lies in it. Similarly every
prime containing \(J+\Delta\mathcal R\) contains \(\mathfrak b\)
or \(\mathfrak m_D\). The Noetherian property now gives a single
integer \(n\) satisfying
\[
  (I_D\mathfrak b)^n\subseteq J,\qquad
  (\mathfrak m_D\mathfrak b)^n\subseteq J+\Delta\mathcal R.
\]
Indeed \(\mathcal R\) has finitely many generators over \(D\) by
the module finiteness proved in Section~\ref{sec:deformation-chart},
and an ideal with finitely many generators in a radical has a
power in the ideal. If all three \(d_i\ge n\), each element of
\(\mathcal R_d\) belongs to \(\mathfrak b^n\), by expanding it
into products of generating sections. Comparing homogeneous
components therefore gives
\[
                     I_D^nP_d=0,\qquad
                  \mathfrak m_D^nP_d\subseteq\Delta P_d .
\]
The ring \(D/I_D^n\) is finite over \(\mathcal O\) because its successive
pieces \(I_D^j/I_D^{j+1}\) are finite over \(D/I_D\). Thus \(P_d\)
is finite over \(\mathcal O\) and so is \(D/\operatorname{ann}P_d\).

For any \(\sigma:D\to k\) in characteristic zero whose localization
\((P_d)_{\ker\sigma}\) is nonzero, the first equality implies that
\(\sigma(I_D)=0\). If \(\sigma(\Delta)=0\), then \(\varpi\in\mathfrak m_D\)
is invertible at \(\ker\sigma\), and the second inclusion gives
\((P_d)_{\ker\sigma}=\Delta(P_d)_{\ker\sigma}\). Nakayama's lemma
contradicts its nonzeroness. Consequently every such \(\sigma\)
has \(\sigma(uv)\ne0\) for some \(u\in A_{-2},v\in A_2\).

This \(\sigma\) extends to a mixed geometric solution of the original
equations. In fact \((A/I_0)_0=D/I_D\) is a direct summand of
\(A/I_0\) as a \(D/I_D\)-module by degree projection.
Consequently \((A/I_0)\otimes_{D,\sigma} k\) contains \(k\) as a
direct summand. This is a nonzero algebra of finite type over \(k\),
hence has a \(k\)-point. The product \(uv\) is nonzero there,
forcing both signs to occur. Its restriction to \(D\) is one of
the restrictions of the finite set \(T\) meeting the mixed orbits,
since torus scaling fixes \(D\).

For such a \(\sigma\) the open-set polynomial description above
also bounds local multiplicities. Invert \(\tau=uv\) and write
the weighted sections as scalar coefficients times the three
polynomial variables. Put \(I_2=(F_i^2,G_i,E_i^2,E'_i,Z_i)\) in
\(A\). In any multidegree with all \(d_i\ge1\), the image of the
relations in that degree contains the ideal generated by those
scalar coefficients and the \(Z_i\) times the degree-\(d\) monomial.
These scalar coefficients generate \((I_2)_\tau\cap D_\tau\)
because every generator of \(I_2\) has even weight and, after
inversion, each even-weight piece is a unit multiple of \(D_\tau\).
It follows that \((P_d)_\tau\) is a quotient of the indicated
even-weight piece of \((A/I_2)_\tau\), after identifying that
piece with \(D_\tau/((I_2)_\tau\cap D_\tau)\) by the weight unit.
The same holds after localization at \(\ker\sigma\). Each weight
piece for the original equations has localized length at most one
by Section~\ref{sec:deformation-chart}. The repeated squaring
argument of Section~\ref{sec:single-character} therefore proves
\[
             \operatorname{length}_{D_{\ker\sigma}}
                 (P_d)_{\ker\sigma}\le 2^{r+1+r'} .
\]

\subsection{Euler characteristic as a sum of valuations}

Let \(P=P_d\) and denote \(m_\sigma=
\operatorname{length}_{D_{\ker\sigma}}P_{\ker\sigma}\).
For \(a\in D\) with \(\sigma(a)\ne0\) on the generic support,
both the kernel and cokernel of \(a:P\to P\) have finite length.
Moreover
\[
  \operatorname{length}_D(P/aP)-\operatorname{length}_D P[a]
       =\sum_{\sigma:D\to k,\;m_\sigma>0}
                    m_\sigma v_{\mathcal O}(\sigma(a)).
\]
To prove the identity remove the \(\mathcal O\)-torsion \(P_{\rm tor}\).
It has finite length, so its kernel and cokernel under an endomorphism
have the same length. The snake sequence shows that their difference
is additive in short exact sequences. On \(P/P_{\rm tor}\), which
is finite free over \(\mathcal O\), multiplication by \(a\) has
nonzero determinant. On the generic fiber all simple factors are
residue fields of the finite-dimensional algebra
\((D/\operatorname{ann}P)\otimes_{\mathcal O}F\), and \(a\) is
nonzero in each occurring field. Smith normal form gives
\(v_{\mathcal O}(\det a)\) for the length of its cokernel. The
\(D\)- and \(\mathcal O\)-lengths of finite-length factors agree
because their local residue fields agree.

In a composition series over the generic algebra each residue
field \(E/F\) occurs as many times as the length of the component
localized at its corresponding maximal ideal. These lengths are
precisely the \(m_\sigma\), since \(\varpi\) is already invertible
at \(\ker\sigma\). The determinant on such a simple factor is
\(\operatorname{Norm}_{E/F}(a)\). The extension is separable in
characteristic zero. Its norm factors as the product over embeddings
\(E\hookrightarrow k\), which are exactly the \(\sigma\)'s with this
kernel. Taking valuations proves the formula, including its
multiplicities.

Each contributing \(\sigma\) restricts from \(T\) and factors
through \(D/I_D\), a finite \(\mathcal O\)-algebra. Thus
\(\sigma(a)\) is integral over \(\mathcal O\) and its valuation
is nonnegative. The one-element Euler characteristic in degree
\(d\) is exactly the kernel-cokernel difference just computed.
We have proved, for all sufficiently large \(d\),
\[
        \chi(a;Q)_d
            \le 2^{r+1+r'} \sum_{\xi\in T}v_{\mathcal O}(\xi(a)).
\]
The argument covers each restriction once before enlarging the
sum to \(T\). The source loci are \(\mathrm{ExistsSecNotMem}\),
\(\mathrm{SameDegreeZeroOfRestriction}\),
\(\mathrm{EventuallyMQGrFiniteAnnihilator}\),
\(\mathrm{LocalizedLengthLeOfSurjective}\) and
\(\mathrm{HerbrandFinLen}\). The proof here uses the stated finite
presentation, point condition, and cotangent spanning hypotheses
from the arithmetic model. The lower bound and middle Euler
comparison in \eqref{eq:single-rees} are separate assertions.
